\section{监测设计杠杆、时间自相关与预算优化}
\label{sec:part09}

\subsection{多期平均误差分解公式的完全推导}

\begin{theorembox}{命题 6 的完全证明（监测设计三大杠杆的误差分工）}
在假设 H1--H5 下，设环境序列自相关函数为 $\gamma_R(h) \equiv \Cov(R_t, R_{t+h})/v_R$，标准化报告冲击自相关为 $\gamma_u(h) \equiv \Cov(u_t, u_{t+h})/\tau^2$。采用 $T$ 期算术均值 $\bar{R}_T = \frac{1}{T}\sum_{t=1}^T \widehat{R}_t$ 估计跨期均值 $\Rbar$ 时，其均方误差严格服从：
\begin{equation}
  \E[(\bar{R}_T - \Rbar)^2] = \frac{1}{T^2}\sum_{s=1}^T \sum_{t=1}^T \Big\{ v_R \gamma_R(|s-t|) + \tau^2 \gamma_u(|s-t|)\big[\Rbar^2 + v_R \gamma_R(|s-t|)\big] \Big\} + \frac{A_H}{Tm}.
  \label{eq:alloc-decomp-general}
\end{equation}
\end{theorembox}

\begin{proof}
将估计误差拆解为抽样扰动与宏观潜在状态扰动：
\begin{equation}
  \bar{R}_T - \Rbar = \frac{1}{T}\sum_{t=1}^T \left( \frac{C_t}{m} - u_t R_t \right) + \frac{1}{T}\sum_{t=1}^T (u_t R_t - \Rbar) \equiv \Xi_1 + \Xi_2.
\end{equation}
\begin{enumerate}
  \item \textbf{正交性证明：}在给定整个潜在过程轨道 $\mathcal{G} \equiv (R_1, \dots, R_T, p_1, \dots, p_T)$ 的条件下，$\E(C_t / m \mid \mathcal{G}) = \frac{p_t n R_t}{\rho n} = u_t R_t$。因此 $\E(\Xi_1 \mid \mathcal{G}) = 0$。由重期望法则：
  \begin{equation}
    \E(\Xi_1 \Xi_2) = \E\big[ \E(\Xi_1 \Xi_2 \mid \mathcal{G}) \big] = \E\big[ \Xi_2 \cdot \E(\Xi_1 \mid \mathcal{G}) \big] = \E[\Xi_2 \cdot 0] = 0.
  \end{equation}
  故两项不相关，总均方误差为各自方差之和：$\E[(\bar{R}_T - \Rbar)^2] = \Var(\Xi_1) + \Var(\Xi_2)$。
  \item \textbf{抽样项 $\Var(\Xi_1)$：}由条件独立性假设 H6，各期的观测抽样在给定潜在状态下条件独立：
  \begin{equation}
    \Var(\Xi_1) = \frac{1}{T^2}\sum_{t=1}^T \E\big[\Var(C_t / m \mid R_t, p_t)\big].
  \end{equation}
  由命题 3 的推导，$\E[\Var(C_t / m \mid R_t, p_t)] = \frac{A_H}{m}$。因此：
  \begin{equation}
    \Var(\Xi_1) = \frac{1}{T^2} \sum_{t=1}^T \frac{A_H}{m} = \frac{T A_H}{T^2 m} = \frac{A_H}{Tm} = \frac{A_H}{B},
  \end{equation}
  其中 $B = Tm$ 是总观测病例量。
  \item \textbf{宏观项 $\Var(\Xi_2)$：}计算双重求和协方差：
  \begin{equation}
    \Var(\Xi_2) = \frac{1}{T^2}\sum_{s=1}^T \sum_{t=1}^T \Cov(u_s R_s, u_t R_t).
  \end{equation}
  利用假设 H5（报告率过程与传播过程相互独立）以及 $\E(u_t) = 1$：
  \begin{align}
    \Cov(u_s R_s, u_t R_t) &= \E(u_s u_t R_s R_t) - [\E(u_t R_t)]^2 \nonumber \\
    &= \E(u_s u_t) \, \E(R_s R_t) - \Rbar^2 \nonumber \\
    &= \big[ 1 + \tau^2 \gamma_u(|s-t|) \big] \big[ \Rbar^2 + v_R \gamma_R(|s-t|) \big] - \Rbar^2 \nonumber \\
    &= \Rbar^2 + v_R \gamma_R(|s-t|) + \tau^2 \gamma_u(|s-t|)\big[\Rbar^2 + v_R \gamma_R(|s-t|)\big] - \Rbar^2 \nonumber \\
    &= v_R \gamma_R(|s-t|) + \tau^2 \gamma_u(|s-t|)\big[\Rbar^2 + v_R \gamma_R(|s-t|)\big].
  \end{align}
\end{enumerate}
将两部分合并，即完全得出式~\eqref{eq:alloc-decomp-general}。
\end{proof}

\subsection{独立环境下的成本优化与 Neyman 费用分配}

当环境与报告在各期相互独立（$\gamma_R(h) = \gamma_u(h) = \mathbb{I}(h=0)$），且报告标准化（$\tau = 0$）时，多期平均方差简化为式~\eqref{eq:alloc-indep}：
\begin{equation}
  V_T(m) = \frac{v_R}{T} + \frac{A}{Tm}.
  \label{eq:alloc-indep}
\end{equation}

\begin{theorembox}{推导 9.1：预算约束下拉格朗日极值证明}
设每增加一个观测时期的固定场地/网络运维成本为 $c_0$，每追踪采集一例病例的边际成本为 $c_1$。总经费预算为 $\mathcal{B}$，预算约束为：
\begin{equation}
  T(c_0 + c_1 m) = \mathcal{B} \implies T = \frac{\mathcal{B}}{c_0 + c_1 m}.
\end{equation}
将 $T$ 代入方差目标函数：
\begin{equation}
  f(m) = \frac{v_R + A/m}{\frac{\mathcal{B}}{c_0 + c_1 m}} = \frac{1}{\mathcal{B}} (c_0 + c_1 m)\left( v_R + \frac{A}{m} \right) = \frac{1}{\mathcal{B}} \left( c_0 v_R + c_1 A + c_1 v_R m + \frac{c_0 A}{m} \right).
\end{equation}
对单期规模 $m$ 求一阶导数：
\begin{equation}
  f'(m) = \frac{1}{\mathcal{B}} \left( c_1 v_R - \frac{c_0 A}{m^2} \right).
\end{equation}
令 $f'(m) = 0$：
\begin{equation}
  c_1 v_R = \frac{c_0 A}{m^2} \iff m^2 = \frac{c_0 A}{c_1 v_R} = \frac{c_0}{c_1} \cdot \frac{A}{v_R} = \frac{c_0}{c_1} \mx.
\end{equation}
由此得出唯一驻点：
\begin{equation}
  \boxed{\, m^* = \sqrt{\frac{c_0}{c_1} \mx} \,}
  \label{eq:opt-m-neyman}
\end{equation}
求二阶导数核验凸性：
\begin{equation}
  f''(m) = \frac{1}{\mathcal{B}} \left( \frac{2 c_0 A}{m^3} \right) > 0, \quad \forall m > 0.
\end{equation}
目标函数在正实轴上严格凸，驻点必为全局唯一最小值。将 $m^*$ 代回目标函数：
\begin{align}
  f(m^*) &= \frac{1}{\mathcal{B}} \left( c_0 v_R + c_1 A + 2\sqrt{c_1 v_R m^* \cdot \frac{c_0 A}{m^*}} \right) \nonumber \\
  &= \frac{1}{\mathcal{B}} \left( c_0 v_R + c_1 A + 2\sqrt{c_0 c_1 v_R A} \right) \nonumber \\
  &= \frac{(\sqrt{c_0 v_R} + \sqrt{c_1 A})^2}{\mathcal{B}}.
\end{align}
\end{theorembox}

\begin{insightbox}{理论联系 9.1：两阶段抽样费用最优分配的同构性}
式~\eqref{eq:opt-m-neyman} 在数学结构上与经典抽样调查理论中的 Neyman--Cochran 两阶段抽样费用最优分配定理 \citep{cochran1977sampling} \textbf{完全同构}。
最优单期样本量取决于固定成本与边际成本之比 $\sqrt{c_0/c_1}$ 以及方差平衡规模 $\sqrt{\mx}$ 的几何均值。
\end{insightbox}

\subsection{时间自相关 AR(1) 环境与关键预算阈值 $\mathcal{B}^\dagger$}

当真实传播条件存在时间记忆性（如天气和季节在相邻周保持连贯），设 $R_t$ 服从平稳高斯 AR(1) 过程：$\Cov(R_s, R_t) = v_R \phi^{|s-t|}$（$\phi \in (0, 1)$）。

\begin{derivationbox}{推导 9.2：自回归环境下的总方差与有效样本量}
计算均值方差项 $\Var(\bar{R}_T) = \frac{v_R}{T^2}\sum_{s=1}^T \sum_{t=1}^T \phi^{|s-t|}$：
将方阵沿主对角线与次对角线求和：
\begin{equation}
  \sum_{s=1}^T \sum_{t=1}^T \phi^{|s-t|} = T + 2\sum_{h=1}^{T-1} (T - h) \phi^h = T \left[ 1 + 2\sum_{h=1}^{T-1} \left(1 - \frac{h}{T}\right)\phi^h \right] \equiv T \kappa_T(\phi).
\end{equation}
利用幂级数求和公式计算内层求和：
\begin{align}
  \sum_{h=1}^{T-1} \phi^h &= \frac{\phi(1 - \phi^{T-1})}{1 - \phi}, \\
  \sum_{h=1}^{T-1} h \phi^h &= \phi \frac{\mathrm{d}}{\mathrm{d}\phi} \sum_{h=1}^{T-1} \phi^h = \frac{\phi - T\phi^T + (T-1)\phi^{T+1}}{(1 - \phi)^2}.
\end{align}
通分相消整理，得出精确闭式：
\begin{equation}
  T^2 \Var(\bar{R}_T) = v_R \left[ \frac{T(1+\phi)}{1-\phi} - \frac{2\phi(1-\phi^T)}{(1-\phi)^2} \right].
\end{equation}
定义单位周期有效相对环境方差因子 $e(T) \equiv \frac{\kappa_T(\phi)}{T}$。当预算约束 $T(c_0 + c_1 m) = \mathcal{B}$ 给出 $m = \frac{\mathcal{B} - c_0 T}{c_1 T}$ 时，目标函数关于连续变量期数 $T$ 表达为：
\begin{equation}
  f_\phi(T) = v_R e(T) + \frac{A c_1}{\mathcal{B} - c_0 T}, \qquad 1 \le T < \frac{\mathcal{B}}{c_0}.
\end{equation}
定义边际检验核函数 $g_T(\phi) \equiv -T^2 e'(T)$：
\begin{equation}
  g_T(\phi) = \frac{1+\phi}{1-\phi} - \frac{2\phi}{(1-\phi)^2} \left[ \frac{2(1-\phi^T)}{T} + \phi^T \ln\phi \right].
\end{equation}
经求导可证：$g_T(\phi)$ 关于 $T$ 处处严格单调递增，且在端点满足：
\begin{equation}
  g_1(\phi) < 1, \qquad \lim_{T \to \infty} g_T(\phi) = \frac{1+\phi}{1-\phi} > 1.
\end{equation}
由介值定理，$g_T(\phi) = 1$ 存在唯一的根 $T^\dagger(\phi)$。
对应地，存在\textbf{临界预算阈值}：
\begin{equation}
  \mathcal{B}^\dagger \equiv T^\dagger (c_0 + c_1 m_0^*), \qquad \text{其中 } m_0^* = \sqrt{\frac{c_0}{c_1}\mx}.
\end{equation}
\begin{itemize}
  \item \textbf{小预算体系 ($\mathcal{B} < \mathcal{B}^\dagger$)：}解满足 $T^* < T^\dagger \implies g_{T^*} < 1 \implies m^* > m_0^*$。
  \textbf{最优策略是减少期数，集中预算做大单期规模！}
  \item \textbf{大预算体系 ($\mathcal{B} > \mathcal{B}^\dagger$)：}解满足 $T^* > T^\dagger \implies g_{T^*} > 1 \implies m^* < m_0^*$。
  \textbf{最优策略是缩减单期规模，将资金用于增加独立观测期数！}
\end{itemize}
\end{derivationbox}

\subsection{空间多辖区协同汇总（Spatial Pooling）}

\begin{theorembox}{注 2 的完全证明（空间池化下的最优规模）}
设协同监测网络涵盖 $J$ 个平行辖区，辖区两两之间环境波动的平均相关系数为 $\bar{c} \in [0, 1]$。总预算为 $\mathcal{B} = J(c_0 + c_1 m)$。则最优单辖区规模为：
\begin{equation}
  m^* = \sqrt{\frac{c_0 A}{c_1 v_R (1 - \bar{c})}}.
  \label{eq:opt-spatial}
\end{equation}
\end{theorembox}

\begin{proof}
根据式~\eqref{eq:pooling}，等权汇总 $J$ 个辖区后的总均方误差为：
\begin{equation}
  \text{MSE}_J(m) = v_R \bar{c} + \frac{v_R(1 - \bar{c})}{J} + \frac{A}{Jm}.
  \label{eq:pooling}
\end{equation}
代入预算约束 $J = \frac{\mathcal{B}}{c_0 + c_1 m}$：
\begin{equation}
  \text{MSE}(m) = v_R \bar{c} + \frac{c_0 + c_1 m}{\mathcal{B}} \left[ v_R(1 - \bar{c}) + \frac{A}{m} \right].
\end{equation}
展开方括号项：
\begin{equation}
  \text{MSE}(m) = v_R \bar{c} + \frac{1}{\mathcal{B}} \left[ c_0 v_R(1-\bar{c}) + c_1 A + c_1 v_R(1-\bar{c})m + \frac{c_0 A}{m} \right].
\end{equation}
对单期单辖区规模 $m$ 求导令其等于零：
\begin{equation}
  c_1 v_R (1 - \bar{c}) - \frac{c_0 A}{m^2} = 0 \implies m^2 = \frac{c_0 A}{c_1 v_R (1 - \bar{c})}.
\end{equation}
开平方即得式~\eqref{eq:opt-spatial}。证毕。
\end{proof}
